\begin{thebibliography}{99}
\raggedright

\bibitem{Behrend1946}
F. A. Behrend, \emph{On sets of integers which contain no three terms in
arithmetical progression}, Proc. Nat. Acad. Sci. U.S.A. \textbf{32} (1946),
331--332. \url{https://doi.org/10.1073/pnas.32.12.331}.

\bibitem{BergelsonHostKra2005}
V. Bergelson, B. Host and B. Kra, \emph{Multiple recurrence and nilsequences},
with an appendix by I. Z. Ruzsa,
Invent. Math. \textbf{160} (2005), no.~2, 261--303.
\href{https://doi.org/10.1007/s00222-004-0428-6}{doi:10.1007/s00222-004-0428-6}.

\bibitem{BergelsonLeibman2007}
V. Bergelson and A. Leibman, \emph{Distribution of values of bounded
generalized polynomials}, Acta Math. \textbf{198} (2007), 155--230.

\bibitem{Bloom2016}
T. F. Bloom, \emph{A quantitative improvement for Roth's theorem on arithmetic
progressions}, J. London Math. Soc. (2) \textbf{93} (2016), no.~3, 643--663.
\url{https://doi.org/10.1112/jlms/jdw010}.

\bibitem{BloomSisask2019}
T. F. Bloom and O. Sisask, \emph{Logarithmic bounds for Roth's theorem via
almost-periodicity}, Discrete Anal. \textbf{2019}, Paper No.~4, 20 pp.
\url{https://doi.org/10.19086/da.7884}.

\bibitem{BloomSisask2020}
T. F. Bloom and O. Sisask, \emph{Breaking the logarithmic barrier in Roth's
theorem on arithmetic progressions}, arXiv:2007.03528 (2020), revised 2021.
\url{https://arxiv.org/abs/2007.03528v2}.

\bibitem{BloomSisaskImprovement2023}
T. F. Bloom and O. Sisask, \emph{An improvement to the Kelley--Meka bounds on
three-term arithmetic progressions}, arXiv:2309.02353 (2023).
\url{https://arxiv.org/abs/2309.02353}.

\bibitem{BloomSisaskKelleyMeka}
T. F. Bloom and O. Sisask, \emph{The Kelley--Meka bounds for sets free of
three-term arithmetic progressions}, Essential Number Theory
\textbf{2} (2023), no.~1, 15--44.
\url{https://doi.org/10.2140/ent.2023.2.15};
revised version \href{https://arxiv.org/abs/2302.07211v3}{arXiv:2302.07211v3}
(May 13, 2025).

\bibitem{Bourgain1999}
J. Bourgain, \emph{On triples in arithmetic progression},
Geom. Funct. Anal. \textbf{9} (1999), no.~5, 968--984.
\url{https://doi.org/10.1007/s000390050105}.

\bibitem{Bourgain2008}
J. Bourgain, \emph{Roth's theorem on progressions revisited},
J. Anal. Math. \textbf{104} (2008), 155--192.
\url{https://doi.org/10.1007/s11854-008-0020-x}.

\bibitem{ConlonFoxZhao2015}
D. Conlon, J. Fox and Y. Zhao, \emph{A relative Szemer\'edi theorem},
Geom. Funct. Anal. \textbf{25} (2015), no.~3, 733--762.
\url{https://doi.org/10.1007/s00039-015-0324-9}.

\bibitem{CrootSisask2010}
E. Croot and O. Sisask, \emph{A probabilistic technique for finding
almost-periods of convolutions},
Geom. Funct. Anal. \textbf{20} (2010), no.~6, 1367--1396.
\href{https://doi.org/10.1007/s00039-010-0101-8}{doi:10.1007/s00039-010-0101-8}.

\bibitem{Erdos1961}
P. Erd\H{o}s, \emph{Some unsolved problems},
Magyar Tud. Akad. Mat. Kutat\'o Int. K\"ozl. \textbf{6} (1961), 221--254.

\bibitem{Erdos1974}
P. Erd\H{o}s, \emph{Problems}, Math. Balkanica \textbf{4} (1974),
203--204, Problem~4.33.6.
\url{https://www.renyi.hu/~p_erdos/1974-27.pdf}.

\bibitem{ErdosTuran1936}
P. Erd\H{o}s and P. Tur\'an, \emph{On some sequences of integers},
J. London Math. Soc. \textbf{11} (1936), 261--264.
\url{https://doi.org/10.1112/jlms/s1-11.4.261}.

\bibitem{Furstenberg1977}
H. Furstenberg, \emph{Ergodic behavior of diagonal measures and a theorem
of Szemer\'edi on arithmetic progressions}, J. Analyse Math. \textbf{31}
(1977), 204--256. \url{https://doi.org/10.1007/BF02813304}.

\bibitem{Gowers1998}
W. T. Gowers, \emph{A new proof of Szemer\'edi's theorem for arithmetic
progressions of length four}, Geom. Funct. Anal. \textbf{8} (1998), 529--551.
\url{https://doi.org/10.1007/s000390050065}.

\bibitem{Gowers2001}
W. T. Gowers, \emph{A new proof of Szemer\'edi's theorem},
Geom. Funct. Anal. \textbf{11} (2001), 465--588.
\url{https://doi.org/10.1007/s00039-001-0332-9}.

\bibitem{GowersHahnBanach2010}
W. T. Gowers, \emph{Decompositions, approximate structure, transference,
and the Hahn--Banach theorem}, Bull. London Math. Soc. \textbf{42}
(2010), no.~4, 573--606.
\url{https://doi.org/10.1112/blms/bdq018}.

\bibitem{GreenTaoU32008}
B. Green and T. Tao, \emph{An inverse theorem for the Gowers $U^3(G)$ norm},
Proc. Edinb. Math. Soc. (2) \textbf{51} (2008), no.~1, 73--153.
\href{https://doi.org/10.1017/S0013091505000325}{doi:10.1017/S0013091505000325}.

\bibitem{GreenTaoPrimes2008}
B. Green and T. Tao, \emph{The primes contain arbitrarily long arithmetic
progressions}, Ann. of Math. (2) \textbf{167} (2008), no.~2, 481--547.
\url{https://doi.org/10.4007/annals.2008.167.481}.

\bibitem{GreenTaoDensity2010}
B. Green and T. Tao, \emph{Yet another proof of Szemer\'edi's theorem},
in \emph{An Irregular Mind: Szemer\'edi is 70},
Bolyai Soc. Math. Stud. \textbf{21}, Springer, 2010, 335--342.
\href{https://doi.org/10.1007/978-3-642-14444-8_8}{doi:10.1007/978-3-642-14444-8\_8}.

\bibitem{GreenTaoOrbits2012}
B. Green and T. Tao, \emph{The quantitative behaviour of polynomial
orbits on nilmanifolds},
Ann. of Math. (2) \textbf{175} (2012), no.~2, 465--540.
\href{https://doi.org/10.4007/annals.2012.175.2.2}{doi:10.4007/annals.2012.175.2.2}.

\bibitem{GreenTao2017}
B. Green and T. Tao, \emph{New bounds for Szemer\'edi's theorem, III:
A polylogarithmic bound for $r_4(N)$}, Mathematika \textbf{63} (2017),
no.~3, 944--1040. \url{https://doi.org/10.1112/S0025579317000316}.

\bibitem{GreenTaoZiegler2012}
B. Green, T. Tao and T. Ziegler,
\emph{An inverse theorem for the Gowers $U^{s+1}[N]$-norm},
Ann. of Math. (2) \textbf{176} (2012), no.~2, 1231--1372.
\href{https://doi.org/10.4007/annals.2012.176.2.11}{doi:10.4007/annals.2012.176.2.11}.

\bibitem{HeathBrown1987}
D. R. Heath-Brown, \emph{Integer sets containing no arithmetic progressions},
J. London Math. Soc. (2) \textbf{35} (1987), 385--394.
\url{https://doi.org/10.1112/jlms/s2-35.3.385}.

\bibitem{HostKra2005}
B. Host and B. Kra, \emph{Nonconventional ergodic averages and nilmanifolds},
Ann. of Math. (2) \textbf{161} (2005), no.~1, 397--488.
\href{https://doi.org/10.4007/annals.2005.161.397}{doi:10.4007/annals.2005.161.397}.

\bibitem{KeevashLifshitzLongMinzer2021}
P. Keevash, N. Lifshitz, E. Long and D. Minzer,
\emph{Global hypercontractivity and its applications},
arXiv:2103.04604 (2021).
\url{https://arxiv.org/abs/2103.04604}.

\bibitem{KellerLifshitzMarcus2026}
N. Keller, N. Lifshitz and O. Marcus, \emph{Sharp hypercontractivity for
global functions}, J. Eur. Math. Soc., published online January 9, 2026.
\url{https://doi.org/10.4171/JEMS/1762};
\href{https://arxiv.org/abs/2307.01356v2}{arXiv:2307.01356v2}.

\bibitem{KelleyMeka2023}
Z. Kelley and R. Meka, \emph{Strong bounds for 3-progressions},
in \emph{2023 IEEE 64th Annual Symposium on Foundations of Computer Science (FOCS)},
IEEE, 2023, 933--973.
\href{https://doi.org/10.1109/FOCS57990.2023.00059}{doi:10.1109/FOCS57990.2023.00059};
revised version \href{https://arxiv.org/abs/2302.05537v6}{arXiv:2302.05537v6} (2024).

\bibitem{Leibman2005}
A. Leibman, \emph{Pointwise convergence of ergodic averages for polynomial
sequences of translations on a nilmanifold},
Ergodic Theory Dynam. Systems \textbf{25} (2005), no.~1, 201--213.
\href{https://doi.org/10.1017/S0143385704000215}{doi:10.1017/S0143385704000215}.

\bibitem{LengEquidistribution2023}
J. Leng, \emph{Efficient equidistribution of nilsequences},
\href{https://arxiv.org/abs/2312.10772v5}{arXiv:2312.10772v5}, 2023, revised 2024.

\bibitem{LengSahSawhneyFive}
J. Leng, A. Sah and M. Sawhney, \emph{Improved bounds for five-term arithmetic
progressions}, Math. Proc. Cambridge Philos. Soc. \textbf{177} (2024),
no.~3, 371--413. \url{https://doi.org/10.1017/S0305004124000264}.

\bibitem{LSSSzemeredi2024}
J. Leng, A. Sah and M. Sawhney, \emph{Improved Bounds for Szemer\'edi's
Theorem}, arXiv:2402.17995 (2024).
\url{https://arxiv.org/abs/2402.17995v2}.

\bibitem{LSSInverse2024}
J. Leng, A. Sah, and M. Sawhney,
\emph{Quasipolynomial bounds on the inverse theorem for the Gowers $U^{s+1}[N]$-norm},
\href{https://arxiv.org/abs/2402.17994v3}{arXiv:2402.17994v3}, 2024.

\bibitem{Manners2018}
F. Manners, \emph{Quantitative bounds in the inverse theorem for the Gowers
$U^{s+1}$-norms over cyclic groups},
\href{https://arxiv.org/abs/1811.00718v2}{arXiv:1811.00718v2},
2018, revised 2024.

\bibitem{OpenAIVanDerWaerden2026}
OpenAI, \emph{Quantitative Superexponential Bounds for van der Waerden Numbers},
OpenAI Math Release preprint
\href{https://github.com/openai/math/blob/main/preprints/Quantitative-Superexponential-Bounds-for-van-der-Waerden-Numbers-September-23-2026/paper.pdf}{OAI:Quantitative-Superexponential-Bounds-for-van-der-Waerden-Numbers-September-23-2026},
2026.

\bibitem{Raghavan2026}
R. Raghavan, \emph{Improved Bounds for 3-Progressions},
arXiv:2603.27045v3 (May 15, 2026).
\url{https://arxiv.org/abs/2603.27045v3}.

\bibitem{Roth1953}
K. F. Roth, \emph{On certain sets of integers},
J. London Math. Soc. \textbf{28} (1953), 104--109.
\url{https://doi.org/10.1112/jlms/s1-28.1.104}.

\bibitem{SandersSumsets2008}
T. Sanders, \emph{Additive structures in sumsets},
Math. Proc. Cambridge Philos. Soc. \textbf{144} (2008), no.~2, 289--316.
\href{https://doi.org/10.1017/S030500410700093X}{doi:10.1017/S030500410700093X}.

\bibitem{Sanders2011}
T. Sanders, \emph{On Roth's theorem on progressions},
Ann. of Math. (2) \textbf{174} (2011), no.~1, 619--636.
\url{https://doi.org/10.4007/annals.2011.174.1.20}.

\bibitem{SandersBogolyubov2012}
T. Sanders, \emph{On the Bogolyubov--Ruzsa lemma},
Anal. PDE \textbf{5} (2012), no.~3, 627--655.
\href{https://doi.org/10.2140/apde.2012.5.627}{doi:10.2140/apde.2012.5.627}.

\bibitem{SchoenSisask2016}
T. Schoen and O. Sisask, \emph{Roth's theorem for four variables and additive structures in sums of sparse sets},
Forum Math. Sigma \textbf{4} (2016), e5.
\href{https://doi.org/10.1017/fms.2016.2}{doi:10.1017/fms.2016.2};
\href{https://arxiv.org/abs/1408.2568}{arXiv:1408.2568}.

\bibitem{Szemeredi1969}
E. Szemer\'edi, \emph{On sets of integers containing no four elements in
arithmetic progression}, Acta Math. Acad. Sci. Hungar. \textbf{20} (1969),
89--104. \url{https://doi.org/10.1007/BF01894569}.

\bibitem{Szemeredi1975}
E. Szemer\'edi, \emph{On sets of integers containing no $k$ elements in arithmetic progression},
Acta Arith. \textbf{27} (1975), 199--245.
\href{https://doi.org/10.4064/aa-27-1-199-245}{doi:10.4064/aa-27-1-199-245}.

\bibitem{Szemeredi1990}
E. Szemer\'edi, \emph{Integer sets containing no arithmetic progressions},
Acta Math. Hungar. \textbf{56} (1990), 155--158.
\url{https://doi.org/10.1007/BF01903717}.

\bibitem{vanDerWaerden1927}
B. L. van der Waerden, \emph{Beweis einer Baudetschen Vermutung},
Nieuw Arch. Wisk. (2) \textbf{15} (1927), 212--216.

\bibitem{Ziegler2007}
T. Ziegler, \emph{Universal characteristic factors and Furstenberg averages},
J. Amer. Math. Soc. \textbf{20} (2007), no.~1, 53--97.
\href{https://doi.org/10.1090/S0894-0347-06-00532-7}{doi:10.1090/S0894-0347-06-00532-7}.

\end{thebibliography}
